% Article: a general research paper.
% Replace the title, authors and body text; keep the structure you need.
% Bibliography: biblatex with Biber, entries in references.bib.
\documentclass[11pt]{article}

% ---- Page and language ------------------------------------------------------
\usepackage[a4paper,margin=2.5cm]{geometry} % letterpaper for US Letter
\usepackage[english]{babel}                 % e.g. [ngerman], [french]
\usepackage{microtype}                      % finer spacing and line breaks

% ---- Mathematics, figures and tables ---------------------------------------
\usepackage{amsmath,amssymb}
\usepackage{graphicx}  % \includegraphics for PNG, JPEG and PDF images
\usepackage{booktabs}  % \toprule, \midrule, \bottomrule
\usepackage{tikz}      % the example figure is drawn with TikZ

% ---- References -------------------------------------------------------------
% style=authoryear gives (Author 2024); style=numeric gives [1].
\usepackage[style=authoryear,backend=biber]{biblatex}
\addbibresource{references.bib}
\usepackage{csquotes}

% Load hyperref late; cleveref after it.
\usepackage[hidelinks]{hyperref}
\usepackage[capitalise]{cleveref}

\title{A Concise Title That Says What You Found}
\author{%
  First Author\thanks{Department, University. \texttt{first@example.org}}
  \and
  Second Author\thanks{Institute, City.}}
\date{\today}

\begin{document}
\maketitle

\begin{abstract}
  One paragraph: the question, what you did, what you found, and why it
  matters. Most readers decide here whether to read on, so keep it to about
  150--250 words.
\end{abstract}

\section{Introduction}
\label{sec:introduction}
State the problem and why it is worth solving. Place your work among earlier
studies \parencite{doe2021,roe2019} and end with your contribution. The rest of
the paper is organised as follows: \cref{sec:method} describes the method and
\cref{sec:results} the results.

\section{Method}
\label{sec:method}
Describe what you did in enough detail that someone else could repeat it.
Number the equations you refer to:
\begin{equation}
  \label{eq:model}
  y_i = \beta_0 + \beta_1 x_i + \varepsilon_i,
  \qquad \varepsilon_i \sim \mathcal{N}(0,\sigma^2).
\end{equation}
\Cref{eq:model} is then estimated by least squares \parencite{roe2019}.

\section{Results}
\label{sec:results}
\Cref{fig:trend} shows the trend and \cref{tab:summary} summarises the data.

\begin{figure}[htbp]
  \centering
  % Replace this drawing with your own image, for example:
  % \includegraphics[width=0.7\linewidth]{figures/trend.png}
  \begin{tikzpicture}[scale=0.9]
    \draw[->] (0,0) -- (6.2,0) node[right] {$x$};
    \draw[->] (0,0) -- (0,4.2) node[above] {$y$};
    \foreach \x/\y in {0.5/0.9,1.2/1.4,1.9/1.6,2.6/2.3,3.3/2.4,4.0/3.1,4.7/3.2,5.4/3.8}
      \fill (\x,\y) circle (2pt);
    \draw[thick,blue] (0.3,0.75) -- (5.8,3.95);
  \end{tikzpicture}
  \caption{Observations and the fitted line from \cref{eq:model}.}
  \label{fig:trend}
\end{figure}

\begin{table}[htbp]
  \centering
  \caption{Summary of the observations.}
  \label{tab:summary}
  \begin{tabular}{lrr}
    \toprule
    Group   & $n$ & Mean \\
    \midrule
    Control & 24  & 1.82 \\
    Treated & 26  & 2.47 \\
    \bottomrule
  \end{tabular}
\end{table}

\section{Discussion}
What do the results mean, what are their limits, and what should come next?
As \textcite{doe2021} put it, \enquote{a clear question is half the answer}.

\section{Conclusion}
Summarise the contribution in a few sentences.

\section*{Acknowledgements}
Thank funders, colleagues and data providers here.

\printbibliography

\end{document}
